
\documentclass[12pt,a4paper]{article}
\textwidth 16cm
\topmargin -2.0cm
\oddsidemargin 0cm
\textheight 25cm
\begin{document}
  \title{On Gr\"obner Bases in polynomial rings\\
over Von Neumann regular rings\\
- Abstract -}
  \author{Yosuke Sato \\
Department of Mathematical Sciences, Ritsumeikan University, Japan, \\
ysato@theory.cs.ritsumei.ac.jp}
  \date{}
  \maketitle
\noindent
A commutative ring $R$ with the following property is called 
a Von Neumann regular ring:
\[ \forall a \in R\: \exists b \in R\: a^2b = a\:. \]
Gr\"{o}bner bases of polynomial rings over such rings are first introduced 
in [9]. At the almost same time, we independently came up with the exactly
same idea in [1].
The purpose of [9] was essentially simultaneous computations of 
Gr\"{o}bner bases of a given finite set of polynomials 
in polynomial rings over different finite fields.
On the other hand, our purpose was to provide Gr\"{o}bner bases technique
for boolean constraint solving. 
After that, our work have been developed in 
[3, 4, 5, 6] and a free software is opened in [2].
Recently we found that we can construct special types of comprehensive
Gr\"{o}bner bases as examples of Gr\"{o}bner bases in polynomial rings
over Von Neumann regular rings([7, 8]).
In my talk, I will give a summary of a series of our work.
\begin{thebibliography}{MMa00a}
\bibitem[1]{1}
        Sakai, ~K.,Sato,~Y. and Menju,~S (1990). Boolean gr\"obner bases(revised).
        Technical Report 613, ICOT, 1990.
\bibitem[2]{2}
        Sato,~Y. (1995,1996,1998). Set Constraint Solvers. A free software developed as
        a Research Funding Program of AITEC, Research Institute For Advanced
        Information Technology, \\
        http://www.icot.or.jp/ARCHIVE/Museum/FUNDING/funding-95-E.html\\
        http://www.icot.or.jp/ARCHIVE/Museum/FUNDING/funding-96-E.html\\
        http://www.icot.or.jp/ARCHIVE/Museum/FUNDING/funding-98-E.html\\
\bibitem[3]{3}
        Sato,~Y. (1996). Application of Groebner basis in constraint of non-numerical domains.
        presented in The 2nd IMACS Conference on Applications of Computer Algebra.
\bibitem[4]{4}
        Sato,~Y. (1997). Set Constraint Solver 
        - Groebner bases for non-numerical domains -.
        International Symposium on Symbolic and Algebraic Computation(ISSAC 97), 
        Poster Abstracts pp 13--14.
 \bibitem[5]{5}
	Sato,~Y. (1998). A new type of canonical Gr\"obner bases
	in polynomial rings over Von Neumann regular rings.
	International Symposium on Symbolic and Algebraic
	Computation (ISSAC~98), Proceedings, 317--321
 \bibitem[6]{6}
        Sato,~Y. and Suzuki,~A.
        (1999). Parallel Computation of Boolean Gr\"{o}bner bases.
        The Fourth Asian Technology Conference in Mathematics(ATCM 99),
        Proceedings pp 265--274.
 \bibitem[7]{7}
        Sato,~Y. and Suzuki,~A.
        (2000). Comprehensive Gr\"{o}bner bases.
        presented in The 6th IMACS Conference on Applications of Computer Algebra.
 \bibitem[8]{8}
        Sato,~Y. and Suzuki,~A.
        (2000). Gr\"{o}bner bases in polynomial rings over Von Neumann regular rings - their applications - (Extended abstract)
        Proceedings of the Fourth Asian Symposium in Computer Mathematics(ASCM 2000),
        Lecture Notes Series on Computimg Vol.8, World Scientific, pp 59--62.
 \bibitem[9]{9}
	Weispfenning,~V. (1989). Gr\"obner bases in polynomial
	ideals over commutative regular rings,
	EUROCAL~'87, J.~H.~Davenport Ed., 
	Springer LNCS{\bf~378}, 336--347.
\end{thebibliography}
\end{document}
